En los centros de distribución, el problema de preparación de pedidos ({\bf OPP}, del inglés {\em Order Picking Problem}) consiste
en asignar lotes de pedidos a distintos operarios, determinar el orden de ejecución y establecer la ruta de recolección en el almacén.
El objetivo es minimizar el tiempo total de preparación, así como las penalidades por adelanto o demora respecto a las fechas de entrega
comprometidas. Aunque la resolución exacta de este problema garantiza la optimidad, el tiempo de cómputo crece rápidamente con el tamaño del
almacén y la cantidad de pedidos. Esto dificulta su aplicación en escenarios operativos reales, donde las decisiones exigen tiempos
de respuesta acotados.

Este trabajo toma como punto de partida un modelo de programación lineal entera mixta de la literatura para el OPP. Tomándolo como guía,
se formula un modelo lineal propio, al que llamamos {\em monolítico}, que incorpora velocidades heterogéneas entre operarios y restricciones
de cadena de frío. Luego, con el fin de reducir sus tiempos de ejecución y abordar instancias de mayor escala, este modelo se optimiza mediante
su descomposición en dos etapas, que se resuelven una después de la otra. Primero se resuelve el ruteo de cada lote y luego, la asignación y la
secuenciación de los lotes. A este segundo modelo lo llamamos {\em Split}. Debido a que el modelo Split sigue resultando demasiado costoso y no se logra resolver en tiempos
prácticos para instancias moderadamente grandes, se desarrollan dos metaheurísticas basadas en {\em Hill Climbing} y {\em Ant Colony Optimization},
equivalentes a los modelos lineales en las características que contemplan.

Los modelos se evalúan sobre un conjunto de escenarios de prueba funcionales y de estrés. A partir de la proporción de instancias resueltas, el
tiempo de ejecución y la calidad de la solución respecto del óptimo, se busca caracterizar el compromiso entre calidad, tiempo de cómputo y robustez,
para identificar bajo qué condiciones resulta conveniente reemplazar el modelo exacto por una heurística, y cuál de las dos metaheurísticas
elegidas es más favorable.
